% =====================================================================
%  LDA-T318 — Lab 2 briefing — Fine-tune vs prompt
%  Thursday 5 November 2026
%  Build:  latexmk -xelatex main.tex
% =====================================================================
\documentclass[aspectratio=169,10pt]{beamer}

% Path has no parent-directory step: Overleaf resolves \input from the project root and
% will not read above it. Locally, the Makefile puts common/ on TEXINPUTS.
% =====================================================================
%  LDA-T318 — Large Language Models and Generative AI for NLP
%  Shared Beamer preamble for all lecture decks.
%  Aarne Talman, University of Helsinki, 2026.
%
%  Engine:  XeLaTeX or LuaLaTeX preferred (gives the Fira font set).
%           pdfLaTeX also works; Metropolis falls back to Latin Modern.
%  Build:   latexmk -xelatex main.tex      (from the deck's own directory)
%
%  Every deck does:
%      \documentclass[aspectratio=169,10pt]{beamer}
%      \input{../common/preamble.tex}
% =====================================================================

% ---------------------------------------------------------------- theme
\usetheme{metropolis}
\metroset{
  titleformat=smallcaps,
  numbering=fraction,
  progressbar=frametitle,
  block=fill,
  sectionpage=progressbar,
  subsectionpage=none,
}

% ---------------------------------------------------------------- packages
% inputenc/fontenc are required under pdfLaTeX and actively harmful under
% XeLaTeX/LuaLaTeX, where Metropolis loads fontspec to get Fira. Guard them.
\usepackage{iftex}
\ifPDFTeX
  \usepackage[utf8]{inputenc}
  \usepackage[T1]{fontenc}
  \usepackage[sfdefault]{FiraSans}
  \usepackage{FiraMono}
\else
  % Metropolis only gets Fira if fontspec can actually find it. TeX Live ships
  % the OTFs but does NOT register them with fontconfig, so a plain
  % \setsansfont{Fira Sans} silently falls back to Latin Modern and the deck
  % stops looking like Metropolis. Load them by filename instead.
  \usepackage{fontspec}
  \setsansfont{FiraSans}[
    Extension      = .otf,
    UprightFont    = *-Book,
    ItalicFont     = *-BookItalic,
    BoldFont       = *-Medium,
    BoldItalicFont = *-MediumItalic]
  \setmonofont{FiraMono}[
    Extension   = .otf,
    UprightFont = *-Regular,
    BoldFont    = *-Medium,
    Scale       = 0.92]
  \renewcommand{\familydefault}{\sfdefault}
\fi
\usepackage[english]{babel}

\usepackage{amsmath,amssymb}
\usepackage{booktabs}
\usepackage{array}
\usepackage{multicol}
\usepackage{xcolor}
\usepackage{graphicx}
\usepackage{tikz}
\usepackage{pgfplots}
\usepackage{listings}
\usepackage{appendixnumberbeamer}

\pgfplotsset{compat=1.17}
\usetikzlibrary{positioning,arrows.meta,fit,backgrounds,calc,matrix,shapes.geometric,decorations.pathreplacing}

% =====================================================================
%  University of Helsinki palette
%  Primary brand colour is UH Hray/black with the "HY blue" accent.
%  Secondary accents chosen for projector contrast, not for print.
% =====================================================================
\definecolor{UHBlue}{HTML}{0C5A9E}   % primary accent
\definecolor{UHDark}{HTML}{15202B}   % body text / dark backgrounds
\definecolor{UHGrey}{HTML}{5A6672}   % secondary text
\definecolor{UHLight}{HTML}{EEF2F6}  % block backgrounds
\definecolor{UHOrange}{HTML}{D1601A} % alerts, gotchas
\definecolor{UHGreen}{HTML}{1F7A5C}  % examples, "this works"
\definecolor{UHRed}{HTML}{A62B2B}    % failure modes
\definecolor{UHSand}{HTML}{C8B88A}   % neutral third accent

\setbeamercolor{normal text}{fg=UHDark, bg=white}
\setbeamercolor{alerted text}{fg=UHOrange}
\setbeamercolor{example text}{fg=UHGreen}
\setbeamercolor{frametitle}{fg=white, bg=UHBlue}
\setbeamercolor{title separator}{fg=UHBlue}
\setbeamercolor{progress bar}{fg=UHBlue}
\setbeamercolor{progress bar in head/foot}{fg=UHBlue}
\setbeamercolor{progress bar in section page}{fg=UHBlue}
\setbeamercolor{structure}{fg=UHBlue}
\setbeamercolor{palette primary}{fg=white, bg=UHBlue}   % section pages, standout frames
\setbeamercolor{block title}{fg=white, bg=UHBlue}
\setbeamercolor{block body}{fg=UHDark, bg=UHLight}
\setbeamercolor{block title alerted}{fg=white, bg=UHOrange}
\setbeamercolor{block body alerted}{fg=UHDark, bg=UHLight}
\setbeamercolor{block title example}{fg=white, bg=UHGreen}
\setbeamercolor{block body example}{fg=UHDark, bg=UHLight}
\setbeamercolor{footnote}{fg=UHGrey}

\setbeamerfont{frametitle}{size=\large, series=\bfseries}
\setbeamerfont{footnote}{size=\tiny}
\setbeamerfont{caption}{size=\scriptsize}

% Frame title bar.
%
% Do NOT reimplement the frametitle template here. Metropolis's own template
% already spans \paperwidth (so `bg=UHBlue' above gives the full-width bar),
% already struts two-line titles without collision, and already draws the
% frametitle progress bar. Reimplementing it means touching internals whose
% names contain `@', which is NOT a letter in a .tex file -- `\@empty' parses
% as `\@' followed by the word "empty" and the build breaks.
%
% To change the bar's height, adjust Metropolis's own padding length instead,
% inside \makeatletter so that `@' is a letter:
%
%   \makeatletter
%   \setlength{\metropolis@frametitle@padding}{2.4ex}  % default 2.2ex
%   \makeatother
%
\makeatletter
\setlength{\metropolis@frametitle@padding}{1.2ex}
\makeatother

% No deck uses \framesubtitle; if one ever does, Metropolis ignores it and the
% subtitle should go in the frame body instead.

% =====================================================================
%  Course metadata — decks override \lecturename and \lecturenumber
% =====================================================================
\newcommand{\coursecode}{LDA-T318}
\newcommand{\coursename}{LLMs and Generative AI for NLP}
\newcommand{\courseteacher}{Aarne Talman}
\newcommand{\courseinstitution}{University of Helsinki}
\providecommand{\lecturename}{}
\providecommand{\lecturedate}{}

% Footline: course code on the left, lecture name centred, n/N on the right.
\setbeamertemplate{frame footer}{\coursecode\quad\textbar\quad\lecturename}

% =====================================================================
%  Semantic macros — use these instead of ad-hoc formatting, so that a
%  change of house style is a single edit here.
% =====================================================================

% --- Sourcing. Every quantitative claim on a slide should carry one. ---
\newcommand{\source}[1]{%
  \vfill{\scriptsize\color{UHGrey}\raggedright Source: #1\par}%
}
\newcommand{\sources}[1]{%
  \vfill{\scriptsize\color{UHGrey}\raggedright Sources: #1\par}%
}
% Inline citation chip, e.g. \cit{Raschka 2026a}
\newcommand{\cit}[1]{{\scriptsize\color{UHGrey}[#1]}}

% --- Counter-intuitive point. Three per lecture; these are exam material. ---
\newenvironment{gotcha}[1][Gotcha]{%
  \begin{alertblock}{#1}%
}{%
  \end{alertblock}%
}

% --- Speculative / unconfirmed claim. Research literacy matters here. ---
\newcommand{\rumour}{{\color{UHOrange}\textbf{[unconfirmed]}}\ }

% --- Key takeaway, for the closing slide of each section. ---
\newenvironment{takeaway}{%
  \begin{block}{Take away}%
}{%
  \end{block}%
}

% --- Live demo marker. ---
\newcommand{\demo}[1]{{\color{UHGreen}\textbf{Demo:}} #1}

% --- Emphasis consistent with the palette. ---
\newcommand{\key}[1]{{\color{UHBlue}\textbf{#1}}}
\newcommand{\warn}[1]{{\color{UHOrange}\textbf{#1}}}
\newcommand{\good}[1]{{\color{UHGreen}\textbf{#1}}}
\newcommand{\bad}[1]{{\color{UHRed}\textbf{#1}}}

% --- Units and quantities used constantly across the decks. ---
\newcommand{\unit}[1]{\,\mathrm{#1}}
\newcommand{\flops}{\mathrm{FLOPs}}
% NB: no \toks macro here --- \toks is a TeX primitive (a token register), so
% \newcommand{\toks} is a hard "already defined" error. No deck used it.
\newcommand{\KV}{\mathrm{KV}}
\newcommand{\softmax}{\operatorname{softmax}}

% =====================================================================
%  Code listings — used for config snippets, Slurm scripts, chat templates.
%  Frames containing these MUST be declared [fragile].
% =====================================================================
\lstdefinestyle{uh}{
  basicstyle=\ttfamily\scriptsize\color{UHDark},
  keywordstyle=\color{UHBlue}\bfseries,
  commentstyle=\color{UHGrey}\itshape,
  stringstyle=\color{UHGreen},
  numberstyle=\tiny\color{UHGrey},
  backgroundcolor=\color{UHLight},
  frame=none,
  xleftmargin=0.6em,
  framexleftmargin=0.6em,
  breaklines=true,
  breakatwhitespace=true,
  showstringspaces=false,
  columns=fullflexible,
  keepspaces=true,
  upquote=true,
  literate={-}{{-}}1 {_}{{\_}}1,
}
\lstset{style=uh}

\lstdefinelanguage{slurm}{
  morekeywords={SBATCH,module,srun,source,export,set,echo},
  morecomment=[l]{\#},
  sensitive=true,
}
\lstdefinelanguage{jinja}{
  morekeywords={for,endfor,if,endif,set,in},
  morecomment=[s]{\{\#}{\#\}},
  sensitive=true,
}

% =====================================================================
%  TikZ styles shared across decks (block diagrams, token grids).
% =====================================================================
\tikzset{
  mod/.style     = {draw=UHBlue, fill=UHLight, rounded corners=2pt,
                    minimum height=7mm, minimum width=22mm,
                    font=\scriptsize, align=center},
  modhi/.style   = {mod, draw=UHOrange, fill=UHOrange!12},
  modmute/.style = {mod, draw=UHGrey, fill=white, text=UHGrey},
  flow/.style    = {-{Latex[length=2mm]}, draw=UHGrey, thick},
  flowhi/.style  = {flow, draw=UHOrange},
  resid/.style   = {-{Latex[length=2mm]}, draw=UHBlue, thick, dashed},
  tokcell/.style = {draw=UHGrey!50, minimum size=4.2mm, font=\tiny, inner sep=0pt},
  tokon/.style   = {tokcell, fill=UHBlue!35},
  tokoff/.style  = {tokcell, fill=UHLight},
  tokmask/.style = {tokcell, fill=UHGrey!15, text=UHGrey!60},
  lbl/.style     = {font=\scriptsize\color{UHGrey}},
}

% =====================================================================
%  Presenter notes.
%
%  Notes are written on every slide. By default they are NOT shown, so
%  `latexmk -xelatex main.tex` produces a clean deck for projection.
%
%  To present with notes on a second screen (pdfpc, Skim, Impressive):
%      uncomment the line below, rebuild to main-notes.pdf
%  To print a notes handout:
%      \setbeameroption{show only notes}
% =====================================================================
% \setbeameroption{show notes on second screen=right}

\setbeamertemplate{note page}{%
  {\scriptsize\color{UHGrey}\coursecode\ \textbar\ \lecturename\par}%
  \vskip0.3em\hrule\vskip0.6em%
  {\small\insertnote}%
}

% ---------------------------------------------------------------- misc
\setbeamertemplate{caption}[numbered]
\setbeamertemplate{itemize items}[circle]
\setlength{\parskip}{0.1em}

% ---------------------------------------------------------------- density
% These decks are text-dense and a 16:9 frame is only 90mm tall, so the
% default vertical rhythm costs several lines per slide. Tightened globally
% here rather than slide by slide.
\setbeamersize{text margin left=0.7cm, text margin right=0.7cm}
\makeatletter
\def\@listi{\leftmargin\leftmargini
  \topsep 1.5\p@ \@plus1\p@ \@minus1\p@
  \parsep 0\p@
  \itemsep 1.5\p@ \@plus1\p@ \@minus1\p@}
\let\@listI\@listi
\makeatother

% The title page is a stack of short paragraphs, so a non-zero \parskip
% pushes it past the bottom of the frame. Neutralise it there only.
\addtobeamertemplate{title page}{\setlength{\parskip}{0pt}}{}

% Tighter lists — the decks are dense.
\setbeamertemplate{itemize/enumerate body begin}{\small}
\setbeamertemplate{itemize/enumerate subbody begin}{\footnotesize}

% Table helpers
\newcolumntype{L}[1]{>{\raggedright\arraybackslash}p{#1}}
\newcolumntype{C}[1]{>{\centering\arraybackslash}p{#1}}
\newcolumntype{R}[1]{>{\raggedleft\arraybackslash}p{#1}}

% Allow slightly overfull boxes without a warning storm
\hfuzz=2pt

% =====================================================================
%  Shared course-map frame.
%  Every lecture opens with this, with the current week highlighted:
%      \coursemapframe{1}   ... \coursemapframe{4}
%  Pass 0 to highlight nothing.
% =====================================================================
% \cmrow{current week}{this row's week}{dates}{Tuesday}{Thursday}{teacher}
% Highlights the row in UH blue bold when the two week numbers match.
% Deliberately avoids \rowcolor so that colortbl is not needed.
% The & and \\ must NOT sit inside the \ifnum branches: TeX scans alignment
% cells while skipping the false branch, so an & there breaks the conditional
% ("Incomplete \ifnum"). Keep the branches down to a macro name and let the
% separators live in the macro bodies, which expand after the test resolves.
\newcommand{\cmhi}[5]{%
  \textcolor{UHBlue}{\textbf{#1}} & \textcolor{UHBlue}{\textbf{#2}} &
  \textcolor{UHBlue}{\textbf{#3}} & \textcolor{UHBlue}{\textbf{#4}} &
  \textcolor{UHBlue}{\textbf{#5}}\\%
}
\newcommand{\cmplain}[5]{#1 & #2 & #3 & #4 & #5\\}
\newcommand{\cmrow}[6]{%
  \ifnum#1=#2\relax\expandafter\cmhi\else\expandafter\cmplain\fi
  {#2}{#3}{#4}{#5}{#6}%
}
\newcommand{\coursemapframe}[1]{%
\begin{frame}{Course overview}
\centering
\small
\setlength{\parskip}{0pt}
\renewcommand{\arraystretch}{1.08}
\begin{tabular}{@{}c L{1.7cm} L{3.1cm} L{3.7cm} L{1.4cm}@{}}
\toprule
\textbf{Wk} & \textbf{Dates} & \textbf{Tuesday} & \textbf{Thursday} & \textbf{Teacher}\\
\midrule
\cmrow{#1}{1}{27 / 29 Oct}{Language models}{Pretraining}{Aarne}
\cmrow{#1}{2}{3 / 5 Nov}{Post-training}{Lab 2: tune vs prompt}{Aarne}
\cmrow{#1}{3}{10 / 12 Nov}{Inference}{Lab 3: serve}{Aarne}
\cmrow{#1}{4}{17 / 19 Nov}{Evaluation}{Lab 4: evaluate}{Jussi}
\cmrow{#1}{5}{24 / 26 Nov}{RAG}{Lab 5: RAG}{Dmitry}
\cmrow{#1}{6}{1 / 3 Dec}{Agents}{Lab 6: agent}{Dmitry}
\cmrow{#1}{7}{8 / 10 Dec}{Course recap}{Special topics (TBC)}{All}
\bottomrule
\end{tabular}

\vspace{1.2em}
{\footnotesize Reusing a task and model across labs is recommended, not required.\\
There are no intermediate submissions; the final report and code repository
determine the grade.}
\note{Fifteen seconds, not more. The point of this slide is orientation, not
content: show them the shape of the course and where today sits inside it.
The final week is deliberately left flexible for a recap and special topics.}
\end{frame}%
}


\renewcommand{\lecturename}{Lab 2 — Fine-tune vs prompt}
\renewcommand{\lecturedate}{5 November 2026}

\title{Lab 2: Fine-tune vs prompt}
\subtitle{Does fine-tuning beat your best prompt?}
\author{Aarne Talman}
\institute{\coursecode\ \textbar\ \coursename\ \textbar\ \courseinstitution}
\date{5 November 2026}

\begin{document}

% =====================================================================
\begin{frame}[plain,shrink=8]
  \titlepage
  \note{Today is a lab rather than a lecture. The whole briefing is ten
  minutes, fifteen at the absolute outside, and then the room works while I
  circulate.

  Open by saying that out loud: \emph{I will talk for ten minutes, then you
  start, and I come to you.} It sets the expectation that nobody has to sit
  through forty minutes of slides before touching a keyboard.

  Before minute 0: check that the GPU reservation is active and that at least
  one student has a \texttt{gpuinteractive} session up. If the reservation has
  not landed, say so immediately and switch the room to pairs.}
\end{frame}

\coursemapframe{2}

% =====================================================================
\section{The question}
% =====================================================================

\begin{frame}{One question}
  \begin{block}{Lab 2}
    \centering\large
    For \emph{your} task, does fine-tuning beat your best prompt?
  \end{block}
  \vspace{0.6em}
  \begin{columns}[T]
    \begin{column}{0.52\textwidth}
      That is the whole lab. Everything today is here to help you answer it
      \key{with numbers you produced yourself}.
      \vspace{0.5em}

      Tuesday's lecture gave you the machinery: assistant-only loss, LoRA,
      catastrophic forgetting, and the decision slide
      (\emph{behaviour \textrightarrow\ fine-tune, knowledge \textrightarrow\ retrieve}).
      Today you test it on your own task.
    \end{column}
    \begin{column}{0.48\textwidth}
      \begin{block}{What week 2 adds to your project}
        \small
        A prompting baseline and a fine-tuned model, \textbf{both scored on the
        same held-out set}.\\[0.4em]
        You may serve the adapter you train today in week 3, or use another
        compatible model. Week 4 examines evaluation in more depth.
      \end{block}
    \end{column}
  \end{columns}
  \vspace{0.4em}
  \small Start with \texttt{lab2-example/README.md}; it contains the baseline,
  LoRA, evaluation and smoke-test commands.
  \note{Thirty seconds. The sentence in the box is the only thing they need to
  remember from the briefing.

  Say explicitly that ``fine-tuning won'' is not the expected answer. For
  several of the tasks, XNLI in particular, a good few-shot prompt on the
  instruct model is hard to beat with 2k examples and a LoRA. A report that
  says ``it did not beat the prompt, here is the number, here is why'' scores
  just as well. Remove the incentive to fake a win now, before anyone starts
  tuning their baseline downwards.}
\end{frame}

\begin{frame}{Material for the final report}
  \centering
  \begin{tabular}{@{}L{4.4cm} C{2.4cm} L{4.2cm}@{}}
    \toprule
    \textbf{Setting} & \textbf{Held-out score} & \textbf{What it is}\\
    \midrule
    Zero-shot & & best of the base or instruct model\\
    \addlinespace[0.2em]
    Best prompt & & choose on dev; report held-out score\\
    \addlinespace[0.2em]
    Fine-tuned (LoRA SFT) & & same held-out set, same decoding\\
    \addlinespace[0.2em]
    Forgetting check & before, after & 50 general-knowledge questions\\
    \bottomrule
  \end{tabular}
  \vspace{0.5em}

  \begin{columns}[T]
    \begin{column}{0.48\textwidth}
      \textbf{Plus one loss curve}, with training \emph{and} evaluation loss
      on the same axes.
    \end{column}
    \begin{column}{0.52\textwidth}
      \textbf{Plus one paragraph:} was it worth it? GPU-hours and human hours
      spent, against points gained. Would you ship it?
    \end{column}
  \end{columns}
  \note{Put this on screen and leave it there for a moment. Everyone should
  photograph it.

  Keep the scope small. Four cells are enough for a clear comparison.

  \textbf{Circulating:} the most common failure is someone who has spent forty
  minutes on data cleaning and has no baseline. Ask them: ``which of these four
  cells can you fill in the next ten minutes?'' and make them do that one.}
\end{frame}

\begin{frame}{Today, minute by minute}
  \centering
  \vspace{0.4em}
  \begin{tikzpicture}[x=1.25mm, y=1mm]
    \fill[UHLight]     (0,0)  rectangle (10,13);
    \fill[UHBlue!16]   (10,0) rectangle (35,13);
    \fill[UHOrange!16] (35,0) rectangle (70,13);
    \fill[UHGreen!16]  (70,0) rectangle (90,13);
    \draw[UHGrey] (0,0) rectangle (90,13);
    \draw[UHGrey] (10,0) -- (10,13);
    \draw[UHGrey] (35,0) -- (35,13);
    \draw[UHGrey] (70,0) -- (70,13);
    \node[font=\scriptsize, align=center] at (5,6.5)
      {set\\up};
    \node[font=\scriptsize, align=center] at (22.5,6.5)
      {\textbf{baseline}\\four prompts, scored once};
    \node[font=\scriptsize, align=center] at (52.5,6.5)
      {\textbf{LoRA SFT}\\launch it, then watch the loss};
    \node[font=\scriptsize, align=center] at (80,6.5)
      {\textbf{measure}\\forgetting, save};
    \foreach \x/\t in {0/0, 10/10, 35/35, 70/70, 90/90}
      \node[lbl, anchor=north] at (\x,-1.5) {\t};
    \node[lbl, anchor=north] at (45,-7) {minutes into the session};
  \end{tikzpicture}
  \vspace{0.8em}

  \raggedright
  \begin{columns}[T]
    \begin{column}{0.5\textwidth}
      \small
      \textbf{This briefing stops at minute 10.}
      If you are not running something by minute 15, put your hand up.
    \end{column}
    \begin{column}{0.5\textwidth}
      \small
      The training run is the long pole, so \key{launch it before you finish
      polishing anything else}. It trains while you write the evaluation code.
    \end{column}
  \end{columns}
  \note{Say the four blocks out loud and then give them the strategy:
  \emph{start the training run early and do everything else while it runs}.
  The students who finish are usually the ones who launched a mediocre run at
  minute 35 and improved it later. People still perfecting their dataset at
  minute 60 tend not to finish at all.

  The reservation window runs to roughly 20:00, so runs can finish after the
  session. Say that out loud, because it lowers the panic.

  \textbf{Circulating:} at minute 35, walk the room and ask everyone whether
  their SFT job is started. One sweep at that point catches most of the people
  who would otherwise run out of time.}
\end{frame}

% =====================================================================
\section{Protocol}
% =====================================================================

\begin{frame}{Run four baselines}
  \centering
  \begin{tabular}{@{}l C{3.0cm} C{3.0cm}@{}}
    \toprule
    & \textbf{Zero-shot} & \textbf{Few-shot (5 examples)}\\
    \midrule
    \textbf{Base model} & & \\[0.6em]
    \textbf{Instruct model} & & \\
    \bottomrule
  \end{tabular}
  \vspace{0.9em}

  \begin{columns}[T]
    \begin{column}{0.52\textwidth}
      Compare all four on the \key{development set}, with the \key{same
      decoding parameters} and the \key{same answer-extraction function}.
      \vspace{0.4em}

      \textbf{Whichever scores highest on dev is \emph{the} baseline.} Report
      its test score only after you have made that choice.
      Do not report the instruct zero-shot number just because it was the
      easiest one to run.
    \end{column}
    \begin{column}{0.48\textwidth}
      \begin{block}{A weak baseline fakes a win}
        \small
        Beating a bad prompt tells you nothing. If you have not tried few-shot
        on the instruct model, you do not have a baseline yet.
      \end{block}
    \end{column}
  \end{columns}
  \note{This is the methodological heart of the lab, so spend two minutes here
  even though it is unglamorous.

  Expect base zero-shot to be poor and possibly degenerate. That is the week~1
  point about base models not knowing when to stop, and it is worth reporting
  rather than hiding.

  \textbf{Circulating:} look for people who have run one baseline and moved on.
  The question to ask: ``which of your four is the best on dev, and how do
  you know?''
  Also look for few-shot examples sampled from the test split; see the next
  slide.}
\end{frame}

\begin{frame}{Decoding settings and the held-out set}
  \begin{columns}[T]
    \begin{column}{0.5\textwidth}
      \textbf{Fix these once, use them everywhere}
      \vspace{0.3em}
      \small
      \begin{tabular}{@{}L{2.8cm} L{2.4cm}@{}}
        \toprule
        temperature & 0 (greedy)\\
        \texttt{max\_new\_tokens} & long enough for your longest gold answer\\
        seed & fixed and recorded\\
        prompt template & identical across all four\\
        answer extraction & one shared function\\
        \bottomrule
      \end{tabular}
      \vspace{0.5em}

      \scriptsize Greedy decoding is the right choice here because the task has
      one correct answer and you want the four runs to be comparable.
      (Week 3: temperature 0 is \emph{not} deterministic once requests are
      batched. Today it is close enough.)
    \end{column}
    \begin{column}{0.5\textwidth}
      \textbf{Three splits}
      \begin{itemize}\small
        \item \textbf{train}: the SFT data, 1--5k examples
        \item \textbf{dev}: where you choose prompts, epochs, learning rate
        \item \textbf{held-out test}: \alert{touched once per reported number}
      \end{itemize}
      \vspace{0.5em}
      Take your few-shot examples from \textbf{train or dev}, never from test.
      \vspace{0.5em}
      \small
      Report the split sizes and the metric in the table caption. ``Accuracy''
      on an unnamed subset is not a number anyone can use.
    \end{column}
  \end{columns}
  \note{The held-out discipline is the part people nod at and then violate. Be
  concrete: if you look at a test score, change your prompt and look again, the
  test set has become a dev set and you can no longer report it as held out.

  Students may use the plan from Lab 1 or choose a different task. If they do
  not yet have a held-out set, tell them to split it now with a fixed seed and
  write the sizes down.

  \textbf{Circulating:} ask to see the line of code that makes the split. If it
  is not in the notebook with a seed, it is not reproducible and week 4 will
  not be able to re-score anything.}
\end{frame}

\begin{frame}[shrink=8]{Three ways to silently ruin a baseline}
  \begin{columns}[T]
    \begin{column}{0.5\textwidth}
      \small
      \textbf{1. Answer extraction.}
      The model answered ``She sold 72 clips.'' Your parser wanted a bare
      integer. Score: 0.
      \textit{Fix:} one extraction function, shared by every run, and ten
      printed generations next to their extracted answers.
      \vspace{0.5em}

      \textbf{2. A chat template on a base model.}
      The base model never saw \texttt{<\textbar im\_start\textbar>} or any
      other control token. Those arrive in post-training.
      \textit{Fix:} plain completion-style prompts for the base model,
      \texttt{apply\_chat\_template} only for the instruct model.
      \vspace{0.5em}

      \textbf{3. Contamination.}
      GSM8K, MNLI and MMLU are in everyone's pretraining mix, so part of your
      ``zero-shot'' score may be recall.
      \textit{Fix:} you cannot decontaminate someone else's pretraining, so
      say so in the report. Week 4 treats this properly.
    \end{column}
    \begin{column}{0.5\textwidth}
      \begin{gotcha}[Check the parser before you believe the number]
        A score of exactly 0.0 is a reason to inspect the extraction code
        before concluding that the model failed.
        \vspace{0.4em}
        Before you believe \emph{any} baseline number, print ten raw
        generations beside their extracted answers and read them.
      \end{gotcha}
      \vspace{0.6em}
      \small
      Parser and template errors often make the baseline too low.
      Contamination can make it too high. Either error invalidates the
      comparison.
    \end{column}
  \end{columns}
  \note{This is the methodological core of the lab. Parser and template errors
  usually depress the baseline, while contamination may inflate it. All three
  make the comparison misleading.

  Give the ten-generations advice as a rule rather than a suggestion.
  ``Print ten and read them'' catches most of the broken results I see in
  reports.

  \textbf{Circulating:} any score of 0.00 or 1.00 is a bug until proven
  otherwise. Ask to see raw generations. Second thing to check: is the base
  model being fed a chat template? Ask them to print the exact prompt string.}
\end{frame}

% =====================================================================
\section{Doing it}
% =====================================================================

\begin{frame}[fragile,shrink=8]{Your data: one conversation per example}
\begin{lstlisting}[language=Python]
# one training example = one conversation, in TRL's "messages" format
{"messages": [
  {"role": "system",    "content": "Solve the problem. End with 'Answer: <n>'."},
  {"role": "user",      "content": "Natalia sold clips to 48 friends ..."},
  {"role": "assistant", "content": "48 + 24 = 72. Answer: 72"}]}

from datasets import load_dataset
train = load_dataset("json", data_files="data/train.jsonl")["train"]
dev   = load_dataset("json", data_files="data/dev.jsonl")["train"]
\end{lstlisting}
  \vspace{0.3em}
  \small
  \begin{itemize}
    \item Let TRL apply \key{the model's own chat template}; templates differ
      between model families.
    \item Compute \key{loss on assistant tokens only}, so the reply is the
      training target while the prompt remains context.
    \item Start with \textbf{1--5k examples}. Make each assistant turn match
      the output format your scorer expects.
    \item Before training, print one rendered example and inspect the exact
      string and loss mask.
  \end{itemize}
  \note{Point at the fourth bullet hardest. Most of the gain students see in
  this lab comes from the model learning to emit the answer in a shape their
  parser can read. That is a real and useful result, and it is also why the
  forgetting check matters.

  Tuesday's lecture had the token-mask picture; this is the same idea with a
  flag on it. The underlying mechanism: TRL builds a label tensor the same shape
  as the input ids and sets the prompt positions to $-100$, which the
  cross-entropy loss ignores. That is all ``assistant-only loss'' means.

  \textbf{Circulating:} ask to see one rendered example, the actual string
  after the chat template has been applied. The runs that produce nonsense are
  nearly always from people who never looked at one.}
\end{frame}

% Maintained code lives in lab2-example; avoid projecting a dense code frame.
\iffalse
\begin{frame}[fragile,shrink=13]{The trainer: bf16 LoRA with TRL}
\begin{lstlisting}[language=Python]
from trl import SFTTrainer, SFTConfig
from peft import LoraConfig

# LoRA: freeze the model, train small adapter matrices beside each linear layer
peft_config = LoraConfig(r=16, lora_alpha=32, lora_dropout=0.05,
                         target_modules="all-linear", task_type="CAUSAL_LM")

args = SFTConfig(output_dir="out/lab2-lora",  # adapter and checkpoints land here
                 bf16=True,                  # native on the GH200
                 assistant_only_loss=True,   # grade the reply only, not the prompt
                 max_length=1024, packing=False,
                 per_device_train_batch_size=8,   # 8 x 2 = 16 sequences per weight update
                 gradient_accumulation_steps=2,
                 learning_rate=2e-4, lr_scheduler_type="cosine",
                 warmup_ratio=0.03, num_train_epochs=2,
                 logging_steps=10, eval_strategy="steps", eval_steps=50)  # eval loss every 50

SFTTrainer(model=MODEL_ID, args=args, train_dataset=train,
           eval_dataset=dev, peft_config=peft_config).train()
\end{lstlisting}
  \vspace{0.2em}
  \small
  \textbf{Success looks like:} a progress bar, a \texttt{loss} number every ten
  steps, and an \texttt{eval\_loss} every fifty.
  Keep \texttt{eval\_strategy} switched on today: the evaluation loss is half
  of your loss curve.
  \source{TRL 1.x (SFT, DPO, GRPO and reward trainers in one library) and PEFT;
  API names current as of the course environment build, Oct 2026.}
  \note{Walk four lines and no more: \texttt{bf16=True},
  \texttt{assistant\_only\_loss=True}, the LoRA config, and
  \texttt{eval\_strategy}. The rest they can read.

  If anyone asks about the effective batch size: 8 sequences per step times 2
  accumulation steps is 16 sequences per optimiser update, on one GPU. Halve
  the first and double the second if they hit out-of-memory; the product, and
  therefore the learning dynamics, stays the same.

  If an API name has drifted in the installed TRL version, trust the notebook
  in the course repo over this slide. Say that once.

  \textbf{Circulating:} check \texttt{assistant\_only\_loss}. Without it the
  model trains on its own prompts and the loss curve looks suspiciously good.}
\end{frame}

\fi
\begin{frame}[shrink=9]{Starting hyperparameters, and the three you may tune}
  \begin{columns}[T]
    \begin{column}{0.52\textwidth}
      \small
      \begin{tabular}{@{}L{2.6cm} L{3.1cm}@{}}
        \toprule
        \textbf{Knob} & \textbf{Start at}\\
        \midrule
        LoRA rank & 16\\
        LoRA alpha & 32 (twice the rank)\\
        LoRA dropout & 0.05\\
        target modules & \texttt{all-linear}\\
        learning rate & 0.0002\\
        schedule & cosine, 3\% warmup\\
        epochs & 2\\
        effective batch & 16 sequences (8 $\times$ 2)\\
        max length & 1024 tokens\\
        precision & bf16\\
        \bottomrule
      \end{tabular}
    \end{column}
    \begin{column}{0.48\textwidth}
      \begin{block}{You may tune exactly three}
        \begin{enumerate}\small
          \item \key{learning rate}
          \item \key{number of epochs}
          \item \key{LoRA rank}, the capacity of the adapter
            (move alpha with it, keeping alpha at twice the rank)
        \end{enumerate}
      \end{block}
      \vspace{0.3em}
      \small
      \textbf{Leave everything else alone.} Change \emph{one} knob at a time and
      log every run in one small table: knob, value, train loss, eval loss,
      held-out score.
      \vspace{0.4em}

      That table is more useful than a better score without a record of what
      changed.
    \end{column}
  \end{columns}
  \vspace{0.3em}
  \source{Common LoRA practice and TRL/PEFT defaults. These are
  \textbf{starting points}, so do not quote any of them in the report as a
  finding.}
  \note{The constraint is there for a reason. Given a free afternoon and ten
  knobs, people random-search and learn nothing. Three knobs, changed one at a
  time and written down, gives you a controlled experiment.

  If someone asks why alpha is twice the rank: the adapter's output is scaled
  by $\alpha/r$, so it is the \emph{ratio} that acts like a learning-rate
  multiplier. Keeping the ratio pinned at 2 means that changing $r$ changes the
  capacity without changing the effective step size, which is what makes it a
  controlled experiment. It is a convention rather than a law. Thirty seconds,
  Tuesday covered it.

  \textbf{Circulating:} the tell-tale of a random search is three notebook
  cells with different numbers and no record of which produced what. Hand them
  a structure: one dict per run, appended to a list.}
\end{frame}

\begin{frame}{When the loss falls and nothing else moves}
  \begin{columns}[T]
    \begin{column}{0.54\textwidth}
      \begin{gotcha}[Lower training loss is not a result]
        Training loss down, evaluation loss flat or rising, held-out score
        unchanged: you memorised the training set.
        \vspace{0.4em}
        With LoRA this happens \emph{fast}, because you have more trainable
        parameters than training tokens.
      \end{gotcha}
      \vspace{0.5em}
      \small
      \textbf{Do the arithmetic for your own run.}
      \texttt{model.print\_trainable\_parameters()} gives you the numerator.
      Your token count is roughly
      (examples) $\times$ (tokens per example).
    \end{column}
    \begin{column}{0.46\textwidth}
      \small
      A plausible shape for a 2B model with rank 16 on all linear layers:
      \vspace{0.3em}
      \begin{tabular}{@{}L{3.3cm} R{1.8cm}@{}}
        \toprule
        trainable parameters & $\sim$10--40\,M\\
        2k examples $\times$ 300 tok & $\sim$0.6\,M tok\\
        \bottomrule
      \end{tabular}
      \vspace{0.5em}

      \textbf{What to do about it}
      \begin{itemize}
        \item plot \key{train and eval loss together}, always
        \item stop at the epoch where eval loss turns up
        \item then reduce epochs, then the learning rate, then $r$, in that order
      \end{itemize}
    \end{column}
  \end{columns}
  \vspace{0.2em}
  \source{Parameter counts here are an order-of-magnitude illustration. Print
  the real number for your own model and report that.}
  \note{Second gotcha of the deck, and the one that shows up in the report
  rather than in the session. Half the submitted loss curves have only the
  training loss on them, and a monotonically falling training curve tells you
  nothing at all.

  The parameter-versus-token comparison is deliberately crude. What matters is
  the ratio, and they can work it out for their own run in one line.

  \textbf{Circulating:} look at the plot before you look at the number. If
  there is one line on it, ask where the other one is.}
\end{frame}

% The batch script is provided in lab2-example/slurm.
\iffalse
\begin{frame}[fragile,shrink=12]{If the run does not fit in the session: use a batch job}
\begin{lstlisting}[language=slurm]
#!/bin/bash
# Every #SBATCH line is a request to the queue, not a command. Comments to Slurm.
#SBATCH --account=project_2021893   # course project billed for GPU hours
#SBATCH --partition=gputest         # which queue. gputest first, always
#SBATCH --nodes=1                   # one machine
#SBATCH --gres=gpu:1                # one GH200 GPU on that machine
#SBATCH --cpus-per-task=16          # Grace CPU cores, for the data loader
#SBATCH --mem=64G                   # host RAM, separate from the GPU's 96 GiB
#SBATCH --time=00:15:00             # killed at this point. gputest caps at 15 min
#SBATCH --output=logs/%x-%j.out     # your printout goes here. mkdir logs first!

module purge                        # start from an empty module environment
module load pytorch                 # the course environment

srun python sft_lora.py --config configs/lab2.yaml   # srun: run it on the node
\end{lstlisting}
  \vspace{0.2em}
  \small
  \texttt{sbatch job.sh} submits it and prints a job id. \quad
  \texttt{squeue --me} shows \texttt{PD} (waiting) then \texttt{R} (running).
  \quad \texttt{seff <jobid>} afterwards tells you what it really used.
  \textbf{Success looks like} your training log appearing in \texttt{logs/}.
  \vspace{0.2em}

  Debug on \texttt{gputest} with \texttt{--max\_steps 5}: if that job finishes,
  the script is correct. Only then switch to \texttt{--partition=gpumedium}
  with a realistic \texttt{--time}.
  \source{Partitions and limits: CSC Roihu documentation
  (\texttt{gputest} 15 min; \texttt{gpumedium} 36 h, 1--4 GPUs).}
  \note{Most people will do the SFT interactively in Jupyter today and that is
  fine. This slide exists for two cases: the student whose run is clearly not
  going to finish, and everyone who attempts bonus B.

  For anyone who has never used a batch scheduler: the mental model is a
  ticket queue. You hand Slurm a shell script plus a shopping list of
  resources, it waits until a machine matching the list is free, runs the
  script there with no terminal attached, and writes everything the script
  prints into the \texttt{--output} file. \texttt{\%x} expands to the job name
  and \texttt{\%j} to the job id. The \texttt{logs/} directory must already
  exist. Slurm will not create it, and a missing directory gives you a silent
  failure and an empty log.

  Repeat the rule for the fourth time this term: \texttt{gputest} before
  \texttt{gpumedium}. Fifteen minutes in the short queue is a cheap way to find
  a typo.

  \textbf{Circulating:} when someone's job sits in PD forever, check the
  account, the partition and whether they asked for more GPUs than the
  reservation has. \texttt{squeue --me --start} gives an estimate.}
\end{frame}

\fi
\begin{frame}[shrink=11]{The forgetting check}
  \begin{columns}[T]
    \begin{column}{0.54\textwidth}
      \textbf{50 general-knowledge questions}, unrelated to your task. An MMLU
      slice is fine.
      \vspace{0.4em}

      Score them \key{twice}: on the model you started from, and on the
      fine-tuned model. Same script, same decoding, same extraction.
      \vspace{0.6em}

      Treat this as a \key{regression test}. Nobody cares whether your 2B model
      knows 42\% or 55\% of MMLU; what you are reporting is the
      \textbf{difference} between before and after.
      \vspace{0.6em}

      A few points down is normal and worth one sentence. Twenty points down
      means the adapter ate the model: fewer epochs, lower learning rate, lower
      rank.
    \end{column}
    \begin{column}{0.46\textwidth}
      \begin{block}{Why this is in the lab at all}
        \small
        Tuesday's claim: \emph{always measure what fine-tuning broke, not only
        what it improved}.
        \vspace{0.4em}
        GLM-5 ends its pipeline with self-distillation precisely to recover
        capabilities weakened by the earlier RL stages. Catastrophic forgetting
        turns up as a line item in production recipes.
      \end{block}
      \vspace{0.3em}
      \begin{block}{50 questions is a blunt instrument}
        \small
        With only 50 right/wrong items, one standard error is about
        \textbf{7 percentage points} when accuracy is near 50\%.
        Report the pair of numbers and do not over-read a 3-point change.
      \end{block}
    \end{column}
  \end{columns}
  \source{MMLU-as-regression-test framing: Raschka, \textit{Understanding the 4
  Main Approaches to LLM Evaluation}. Standard error from the binomial, $n=50$.}
  \note{Make the regression-test framing explicit, because students instinctively
  read any benchmark number as a quality score and then feel bad about it.

  The error-bar point is cheap and transferable: it is the first time in the
  course we attach uncertainty to a benchmark number, and week 4 will push on it
  much harder. Two sentences, no more.

  \textbf{Circulating:} the common mistake is scoring the forgetting check only
  after fine-tuning, because the before-number ``seemed obvious''. Without the
  before-number there is no result at all. Ask for it early, before the adapter
  exists.}
\end{frame}

% =====================================================================
\section{Bonus}
% =====================================================================

% Optional extensions remain in the source and README.
\iffalse
\begin{frame}{Bonus A: DPO, optional}
  \alert{Only after the four cells are filled in.}
  \begin{columns}[T]
    \begin{column}{0.52\textwidth}
      \textbf{What to run}
      \begin{itemize}\small
        \item \texttt{DPOTrainer}, starting from \key{your SFT adapter} rather
          than from the base model
        \item preference pairs: \texttt{prompt}, \texttt{chosen},
          \texttt{rejected}
        \item \texttt{beta=0.1} and one epoch
        \item learning rate \texttt{5e-6}, about forty times smaller than SFT
      \end{itemize}
      \vspace{0.3em}
      \textbf{Where pairs come from.} A public preference set, or your own
      verifier labelling completions you sampled.
    \end{column}
    \begin{column}{0.48\textwidth}
      \begin{block}{What to report}
        \small
        The same held-out score and the same forgetting check, as two more rows
        under your table. Nothing else.
      \end{block}
      \vspace{0.3em}
      \small
      \textbf{Expect the known failure mode} from Tuesday: the likelihood of
      the chosen \emph{and} the rejected responses can both fall.
      \vspace{0.3em}

      If the held-out score does not move, that is a finding. Report it with
      the beta you used.
    \end{column}
  \end{columns}
  \source{DPO: Rafailov et al. 2023. Beta and learning rate are
  conventional starting values, not tuned for your task.}
  \note{Thirty seconds in the briefing. Say the gate out loud, four cells
  first, and move on. Anyone who gets here is doing well.

  Beta is the leash: it sets how far DPO is allowed to drag the model away from
  the SFT adapter it started from. If pressed for the detail, DPO's loss is a
  log-sigmoid of $\beta$ times the difference in log-probability ratios between
  the chosen and the rejected response, each ratio taken against the frozen
  reference model. Small $\beta$ means the implicit KL constraint is loose and
  the model may wander far from the reference; large $\beta$ pins it close.
  $0.1$ is the value in the original paper.

  On the failure mode: DPO only ever acts on the gap between the two
  likelihoods, so both can drift downwards together while the gap still widens.

  Spell out where the pairs come from. A public preference set works, and so
  does your own verifier if the task is verifiable: sample a handful of
  completions per training prompt, label the correct ones \texttt{chosen} and
  the rest \texttt{rejected}. That is cheap and it is on-policy. It is also the
  bridge to bonus B, since once you can label your own samples automatically
  you are one step from RLVR.

  \textbf{Circulating:} if someone starts DPO from the base model instead of
  their SFT adapter, stop them; they are measuring a different experiment and it
  will not sit in the same table.}
\end{frame}

\begin{frame}[shrink=14]{Bonus B: GRPO, optional and batch-only}
  \begin{columns}[T]
    \begin{column}{0.54\textwidth}
      \textbf{Natural fit: GSM8K.} The reward is a program, so nothing has to
      be trained to judge the answers:
      \vspace{0.3em}

      \key{reward = (is the final answer right?) + (small bonus for the
      required output shape)}
      \vspace{0.3em}

      \emph{In words:} score each sampled answer 1 or 0 by exact match, add a
      little for formatting it correctly. That is the R1-Zero recipe at toy
      scale.
      \vspace{0.4em}

      \textbf{Settings to start from}
      \begin{itemize}\small
        \item \texttt{GRPOTrainer}, group size 8, so eight sampled answers per
          prompt get compared against each other
        \item \textbf{about 100 steps}, which is the whole bonus
        \item start from your SFT adapter
        \item partition \texttt{gpumedium}, ask for 2 hours
      \end{itemize}
    \end{column}
    \begin{column}{0.46\textwidth}
      \begin{alertblock}{Do not run this in your Jupyter session}
        \small
        Every GRPO step generates eight completions for every prompt in the
        batch \emph{before} it does any learning. Almost all the wall-clock
        time is sampling, and it is \emph{slow}.
        \vspace{0.4em}
        Submit it with \texttt{sbatch}. Debug it on \texttt{gputest} with
        \texttt{--max\_steps 2} first.
      \end{alertblock}
      \vspace{0.4em}
      \small
      \textbf{Expectation management.} One hundred steps will not give you a
      reasoning model. Report the \key{reward curve} and what you learned about
      the machinery. Nobody is expecting a win here.
    \end{column}
  \end{columns}
  \source{GRPO and the accuracy+format reward: DeepSeek-R1 (2025).
  Group size, step count and the two-hour request are starting points for this
  lab, not measured requirements.}
  \note{The reward written symbolically, if anyone wants it:
  $R = R_{\text{accuracy}} + R_{\text{format}}$, with $R_{\text{accuracy}}$ an
  exact-match indicator on the final answer and $R_{\text{format}}$ a small
  constant for emitting the required delimiters. GRPO then normalises the $G$
  rewards within each group to get an advantage per completion. There is no
  learned value model, which is the point of the method and why the group size
  $G$ matters.

  Tuesday explained why GRPO is a bonus and not the main exercise: the
  engineering around it --- zero-variance group filtering, token-level loss,
  clip-higher --- is what decides whether a run survives, and that is a project
  in itself.

  Be blunt about the runtime. The failure mode is a student burning the last
  thirty minutes of a \texttt{gpuinteractive} session on a run that would need
  hours, and losing the main deliverable.

  \textbf{Circulating:} if anyone is running GRPO interactively before their
  table is finished, intervene. This is the one place in the lab where I stop
  someone rather than ask a question.}
\end{frame}

% =====================================================================
\fi
\section{Before you leave}
% =====================================================================

\begin{frame}{The report table and files to keep}
  \begin{columns}[T]
    \begin{column}{0.56\textwidth}
      \small
      \begin{tabular}{@{}L{2.9cm} C{1.3cm} L{2.2cm}@{}}
        \toprule
        \textbf{Setting} & \textbf{Score} & \textbf{Config}\\
        \midrule
        Zero-shot, instruct & & greedy\\
        Few-shot 5, instruct & & greedy\\
        Zero-shot, base & & greedy\\
        Few-shot 5, base & & greedy\\
        \midrule
        \textbf{LoRA SFT} & & rank 16, 2 epochs\\
        \midrule
        Forgetting, before & & 50 Q\\
        Forgetting, after & & 50 Q\\
        \bottomrule
      \end{tabular}
      \vspace{0.4em}

      \scriptsize Caption it with: dataset, split sizes, metric, model IDs.
    \end{column}
    \begin{column}{0.44\textwidth}
      \begin{block}{Keep for the final submission}
        \small
        \begin{itemize}
          \item the notebook, run top to bottom
          \item \texttt{sft\_lora.py} and the Slurm script
          \item \texttt{results/lab2.json} with every number above
          \item the loss curve (train \emph{and} eval)
          \item the LoRA adapter, or its path on Roihu if it is large
        \end{itemize}
      \end{block}
      \vspace{0.2em}
      \small
      There is no intermediate submission. Keep the adapter if you want to
      serve it in week 3, and retain the results for the final report.
    \end{column}
  \end{columns}
  \note{The table has seven rows here and four in the deliverable slide. Say
  why: the four numbers tell the \emph{story} and the full grid is the
  \emph{evidence} behind it. Both go in the report.

  Push the naming point. Clear model and result names save time if students
  reuse this adapter in later labs.

  \textbf{Circulating, last ten minutes:} ask whether results and
  configurations have been saved. A number that only exists in a notebook
  kernel is no use to anyone.}
\end{frame}

\begin{frame}[shrink=8]{Roihu reminders}
  \begin{columns}[T]
    \begin{column}{0.5\textwidth}
      \begin{alertblock}{You lost two days of compute this week}
        \small
        Roihu was down for scheduled maintenance on \textbf{3--4 November}, so
        there was no homework compute between Tuesday's lecture and now.
        \vspace{0.3em}
        That is why the baseline block is at the front of today's session rather
        than assumed as homework.
      \end{alertblock}
      \vspace{0.5em}
      \small
      The lab reservation covers roughly \textbf{16:00--20:00}, so runs started
      in the session can finish after it. Next maintenance: 1--2 December.
    \end{column}
    \begin{column}{0.5\textwidth}
      \textbf{Etiquette, unchanged all term}
      \begin{itemize}\small
        \item \texttt{gputest} (15 min, short queue) before \texttt{gpumedium}
        \item \textbf{close idle GPU Jupyter sessions.} One forgotten session
          holds a whole GH200 for twelve hours, and we share one allocation
        \item never train on a login node; it has no GPU and its CPUs are
          shared
        \item \texttt{seff <jobid>} after every batch job, to see what you
          actually used
      \end{itemize}
      \vspace{0.4em}
      \small
      \textbf{Environment:} build and test the course environment on Roihu's
      ARM login nodes before requesting a GPU. The training configuration uses
      bf16, which the GH200 supports natively.
    \end{column}
  \end{columns}
  \source{CSC Roihu documentation: system description, partitions, course
  maintenance windows.}
  \note{Thirty seconds. The part worth some feeling is the idle sessions.
  Twenty-five people share one allocation, and one forgotten session is a
  twelve-hour GH200 hold that someone else needed.

  Acknowledge the maintenance outage honestly rather than pretending the
  homework was possible. It buys goodwill and it justifies the session plan.

  \textbf{At the end of the session:} say ``close your sessions'' out loud as
  people leave, and watch the queue afterwards.}
\end{frame}

\begin{frame}{Leave this on the wall}
  \centering
  \begin{tikzpicture}[x=1.15mm, y=1mm]
    \fill[UHLight]     (0,0)  rectangle (10,10);
    \fill[UHBlue!16]   (10,0) rectangle (35,10);
    \fill[UHOrange!16] (35,0) rectangle (70,10);
    \fill[UHGreen!16]  (70,0) rectangle (90,10);
    \draw[UHGrey] (0,0) rectangle (90,10);
    \draw[UHGrey] (10,0) -- (10,10); \draw[UHGrey] (35,0) -- (35,10);
    \draw[UHGrey] (70,0) -- (70,10);
    \node[font=\scriptsize, align=center] at (5,5)    {set up};
    \node[font=\scriptsize, align=center] at (22.5,5) {\textbf{baseline}};
    \node[font=\scriptsize, align=center] at (52.5,5) {\textbf{LoRA SFT}};
    \node[font=\scriptsize, align=center] at (80,5)   {\textbf{measure, save}};
    \foreach \x/\t in {0/0, 10/10, 35/35, 70/70, 90/90}
      \node[lbl, anchor=north] at (\x,-1.2) {\t};
  \end{tikzpicture}
  \vspace{0.9em}

  \raggedright
  \begin{columns}[T]
    \begin{column}{0.5\textwidth}
      \small
      \texttt{[\,]}\ four baselines run, same decoding

      \texttt{[\,]}\ best on dev named as \emph{the} baseline

      \texttt{[\,]}\ ten raw generations read with your own eyes

      \texttt{[\,]}\ forgetting check scored \textbf{before} training
    \end{column}
    \begin{column}{0.5\textwidth}
      \small
      \texttt{[\,]}\ LoRA SFT run, train \emph{and} eval loss plotted

      \texttt{[\,]}\ fine-tuned model scored on the held-out set

      \texttt{[\,]}\ forgetting check scored after

      \texttt{[\,]}\ scripts, settings and results saved
    \end{column}
  \end{columns}
  \vspace{0.6em}
  \centering\small
  Stuck for more than ten minutes? Put your hand up. \quad
  Appendix slide has the common errors.
  \note{This is the wall slide. \textbf{Leave it projected for the rest of the
  session.} Do not advance past it, and do not go back to the code slides
  unless someone asks. Everything after this frame is appendix.

  It does two jobs: it answers ``what am I supposed to be doing right now''
  without anyone having to ask, and it gives me a shared vocabulary while
  circulating, so I can point at a line instead of re-explaining it.

  The ten-minute rule is worth saying as I step away from the lectern. In a
  room of twenty-five, the people who suffer silently are the ones who do not
  finish.}
\end{frame}

\appendix

\begin{frame}[shrink=14]{Appendix: things that will go wrong}
  \begin{columns}[T]
    \begin{column}{0.5\textwidth}
      \small
      \textbf{\texttt{CUDA out of memory}}\\
      Halve \texttt{per\_device\_train\_batch\_size}, double
      \texttt{gradient\_accumulation\_steps}, cut \texttt{max\_length}.
      Then \textbf{restart the kernel}, because a dead run keeps its memory.
      \vspace{0.4em}

      \textbf{Score is exactly 0.00 or 1.00}\\
      Answer extraction. Print ten raw generations.
      \vspace{0.4em}

      \textbf{Generations are empty or truncated}\\
      \texttt{max\_new\_tokens} too small, or you forgot
      \texttt{add\_generation\_prompt=True}.
      \vspace{0.4em}

      \textbf{``Tokenizer has no chat template''}\\
      You loaded the \emph{base} model. Either prompt it completion-style, or
      use the instruct tokenizer, and say which in the report.
    \end{column}
    \begin{column}{0.5\textwidth}
      \small
      \textbf{\texttt{pad\_token} is \texttt{None}}\\
      \texttt{tok.pad\_token = tok.eos\_token}, then re-create the trainer.
      \vspace{0.4em}

      \textbf{Loss is \texttt{nan}}\\
      Learning rate too high, or something is running in fp16. We are bf16
      everywhere.
      \vspace{0.4em}

      \textbf{Job stuck in \texttt{PD} (pending) forever}\\
      Wrong account or partition, or more GPUs than the reservation holds.
      \texttt{squeue --me --start} for an estimate, \texttt{scancel} to clear.
      An empty log usually means \texttt{logs/} does not exist.
      \vspace{0.4em}

      \textbf{Training loss suspiciously low from step 1}\\
      \texttt{assistant\_only\_loss} is off. You are training on your own prompts.
    \end{column}
  \end{columns}
  \note{Do not present this frame. It exists so that I can jump to it on the
  projector when the same error appears at three desks in a row, which happens
  in every lab, usually OOM and the missing chat template.

  If a problem here hits more than four people, the course notebook is at
  fault. Write it down during the session and fix the repository while the
  details are still fresh.}
\end{frame}

\end{document}
